\section{Обсуждение, перспективы развития и выводы}

\subsection{Что еще можно было бы попробовать?}

В ходе исследования сформировался ряд перспективных идей, которые могут быть реализованы в качестве дальнейшего развития проекта:
\begin{enumerate}
    \item \textbf{Динамическая кластеризация ориентиров (Voronoi Landmark Pruning)}: Сейчас ориентиры сэмплируются равномерно по буферу. Использование диаграмм Вороного и фильтрации по топологической кривизне позволит сократить число вершин с $N=2000$ до $N=300$ без потери разрешающей способности, ускоряя построение графа в 40 раз.
    \item \textbf{Адаптивный температурный контроллер выхода из стагнации}: Вместо дискретного счетчика застреваний можно внедрить дифференцируемый блок оценки кинематического затухания $\Delta v = \|v_t - \bar{v}\|$, плавно модулирующий энтропию действий $\tau_t = \sigma(\alpha \cdot (v_{\text{crit}} - v_t))$ при контакте со стеной.
    \item \textbf{Сквозная дистилляция траекторий в замкнутом цикле (Closed-Loop Rollout Distillation)}: Обучение трансформера последовательностей не на статических срезах $(s, \mathcal{W})$, а через сквозное дифференцирование развернутых траекторий во времени (BPTT) с использованием дифференцируемых физических движков (Brax / MuJoCo XLA).
    \item \textbf{Ансамблирование окон внимания для сверхбольших лабиринтов (AntMaze Giant)}: Комбинация короткого окна ($W=4$) для прохождения узких створов и длинного окна ($W=16$) для открытых залов.
\end{enumerate}

\subsection{Общие выводы}

В данной работе был пройден полный путь исследования и разработки методов многошагового планирования поверх замороженных Forward-Backward представлений в задаче навигации сложного многосуставного робота:
\begin{itemize}
    \item Показано, почему наивные методы деления пути пополам (Recursive Bisection) ломаются из-за топологической невыпуклости лабиринтов и диффузного просачивания оценок достижимости сквозь стены.
    \item Реализован надежный топологический планировщик Search on the Buffer с квази-метрической стоимостью ребер и непрерывным скользящим упреждением (Sliding Lookahead), решивший проблему застревания на базовом уровне.
    \item Предложена и обучена архитектура амортизированной нейросетевой дистилляции в $O(1)$ (\texttt{Distilled JAX Gated-Attention}) с 4-компонентным физико-ориентированным лоссом, обеспечившая ускорение инференса в 5 раз при росте успешности до **86.0\%** на лабиринте Medium.
    \item Разработан SOTA-метод \textbf{Enhanced Sequence Waypoint Attention Transformer}, объединивший топологический граф Дейкстры, затухание внимания ALiBi, токены кривизны и локально-глобальное окно упреждения, что обеспечило уверенную победу над бейзлайном на лабиринте Large (**66.4\% против 24.0\%**).
\end{itemize}

Полученные результаты наглядно доказывают, что правильное рассуждение о последовательностях интенций в сочетании с физико-согласованными механизмами внимания позволяет эффективно решать длинногоризонтные задачи zero-shot навигации без привлечения вычислительно тяжелых генеративных моделей мира.
